% ECONOMICS_PROBLEM: {"id": "C73-496", "title": "Existence of unbeatable strategies in parallel repeated games", "jel_code": "C73", "source_id": "OP-0169"}
% !TeX program = xelatex
\documentclass[11pt,a4paper]{article}
\usepackage[margin=23mm]{geometry}
\usepackage{fontspec}
\setmainfont{Latin Modern Roman}
\usepackage{amsmath,amssymb,mathtools}
\usepackage{xcolor,enumitem,booktabs,longtable,array,xurl,hyperref}
\definecolor{muted}{HTML}{53636C}
\hypersetup{colorlinks=true,urlcolor=blue,linkcolor=blue}
\setlength{\parindent}{0pt}
\setlength{\parskip}{5pt}
\newcommand{\R}{\mathbb R}
\newcommand{\E}{\mathbb E}
\newcommand{\N}{\mathbb N}
\newcommand{\ind}{\mathbf 1}
\newcommand{\Var}{\operatorname{Var}}
\newcommand{\Cov}{\operatorname{Cov}}
\newcommand{\supp}{\operatorname{supp}}
\DeclareMathOperator*{\argmin}{arg\,min}
\DeclareMathOperator*{\argmax}{arg\,max}
\newcommand{\entrymeta}[3]{{\sffamily\small\color{muted}\textbf{\texttt{#1}}\quad|\quad #2\par #3\par}}
\newcommand{\Problemref}[1]{\texttt{#1}}
\newcommand{\JELref}[2]{\texttt{#2} (JEL \texttt{#1})}
\begin{document}
\section*{Existence of unbeatable strategies in parallel repeated games}
\textbf{Problem ID:} \texttt{C73-496}\quad \textbf{JEL:} \texttt{C73}\par
% BEGIN ECONOMICS STATEMENT
\label{problem:OP-0169}
% GAME_THEORY_PROBLEM: {"local_id": "GT-039", "original_number": 39, "kind": "R", "entry_kind": "statement_only", "published": false, "review_required": true, "category": "game_theory"}
\label{game:GT-039}
{\sffamily\small\color{muted}
{\sffamily\small\color{muted}\textbf{GT-039} | Repeated games and strategy structure\par R | Explicit memory-one specialization of a source characterization agenda\par}
\par}
\subsection*{Definitions and mathematical statement}
Let $G$ be a finite $n$-player stage game, with action sets $A_j$ and payoffs $g_j:A=\prod_jA_j\to\mathbb R$. Play two identical copies simultaneously, indexed by $\mu\in\{1,2\}$, observing the complete sequence of joint profiles $(a_1^t,a_2^t)\in A\times A$. Fix a role $j$. Player $(1,j)$ uses a memory-one kernel $\kappa:A\times A\to\Delta(A_j)$; every other player may use an arbitrary behavioral strategy of the full public history. Specify an arbitrary initial pair $a^0=(a_1^0,a_2^0)$.

Whenever both expected Ces\`aro payoff limits exist, write
\[
 S_{\mu,j}=\lim_{T\to\infty}\frac1T\sum_{t=1}^{T}\mathbb E[g_j(a_\mu^t)].
\]
Call $\kappa$ unbeatable, in this limit-based sense, if for every initial pair and every profile of the other players for which these two limits exist, $S_{1,j}\geq S_{2,j}$. Characterize the finite games and designated roles for which
\[
 \exists\kappa\quad\forall a^0\in A\times A\quad\forall\eta_{-(1,j)},
 \qquad (S_{1,j},S_{2,j}\text{ exist})\Longrightarrow S_{1,j}\geq S_{2,j}.
\]
Thus the comparator is the same role in the other copy, rather than an opponent's payoff in one copy.
\subsection*{Short English statement}
When can a memory-one player guarantee at least the long-run payoff of its counterpart in a parallel copy of the game?
\subsection*{Variants and scope boundaries}
Against nonstationary opponents the displayed limits need not exist. A robust alternative requires $\liminf_T T^{-1}\sum_{t\leq T}\mathbb E[g_j(a_1^t)-g_j(a_2^t)]\geq0$ for all opponents. Bounding the cumulative payoff disadvantage by a constant is stronger still. Memory length and observation of both copies are part of the model; changing either changes the characterization question.
\subsection*{Source relationship and status}
[S1], Definition 1, compares the same player role across parallel copies. Its Section 4 retains general necessary-and-sufficient conditions as a question; Appendix A distinguishes a stronger cumulative-payoff notion. The fixed memory-one characterization above is an editorial specialization of that agenda.
\subsection*{References}
\begingroup\footnotesize\raggedright
\textbf{[S1]} Masahiko Ueda, \emph{Unbeatable Imitation of a Friend}, arXiv:2507.16221v2 (2026 revision of a 2025 preprint). Locators: Section 2, Definition 1, p. 5; Section 4, p. 16 (characterization agenda); Appendix A, Definition 9 (stronger cumulative notion). \url{https://arxiv.org/pdf/2507.16221}
\par\endgroup

\subsection*{Common conventions: Source mathematical conventions}

\textbf{Finite games.} For a finite set $A$, $\Delta(A)$ is its probability
simplex. Players choose independent mixed strategies in Nash equilibrium;
correlation is allowed only when explicitly part of the solution concept.
In a game with payoff functions $u_i$ and strategy profile $x$, an additive
$\varepsilon$ Nash equilibrium satisfies
\[
 u_i(x)\geq u_i(x_i',x_{-i})-\varepsilon
 \quad\text{for every player $i$ and every admissible deviation $x_i'$}.
\]
For numerical approximation, payoffs are normalized as locally specified.
An $\varepsilon$ payoff residual is distinct from distance at most
$\varepsilon$ to the set of exact equilibria. Point convergence, convergence
of empirical play, and convergence in distance to an equilibrium set are
also distinct.

\textbf{Computational representations.} Unless stated otherwise, a complexity
question takes rational data with binary encoding; polynomial time is measured
in total bit length. A fixed-accuracy polynomial algorithm is not necessarily
polynomial in $1/\varepsilon$, and neither is necessarily polynomial in
$\log(1/\varepsilon)$. Succinct games and value, demand, or sampling oracles
must declare their interfaces. Exact real solutions require an explicit
representation; an unspecified list of arbitrary real numbers is not a
finite computational output. A claimed algorithmic barrier is meaningful
only for its specified model and reductions.

\textbf{Dynamic games.} Histories, information, observation, and allowable
randomization are part of the model. A stationary strategy depends only on
the current state; a positional strategy in a deterministic graph game
chooses one action at each controlled vertex. Nash equilibrium at an initial
state and equilibrium after every history are different requirements.
An expectation of a pathwise $\liminf$ payoff, a $\liminf$ of expected averages,
a limiting expected average, and a uniform finite-horizon equilibrium are
different criteria. None is silently substituted for another.

\textbf{Allocation and mechanisms.} A complete allocation assigns every item
exactly once unless otherwise stated. Zero-valued goods, negative values,
capacity constraints, feasibility, lotteries, and copies of goods affect
fairness definitions and are specified locally. Dominant-strategy incentive
compatibility, Bayesian incentive compatibility, universal truthfulness,
and truthfulness in expectation impose different inequalities. Whenever
an objective ratio has zero denominator, its convention is stated or those
instances are excluded.

\textbf{Infinite-dimensional models.} Expectations, integrals, stochastic
equations, and optimization objectives require their local regularity,
integrability, filtrations, and admissible domains. A backward--forward
mean-field system is a boundary-value problem; it is not an initial-value
semiflow without an additional construction. Long-time, large-population,
and vanishing-noise limits require an order of limits and a convergence
topology. If a minimum need not be attained, the objective is its infimum
or supremum and the approximation requirement is stated separately.
% END ECONOMICS STATEMENT
\end{document}
